% =============================================================================
%  Sección 7: Conclusiones
% =============================================================================

\section{Conclusiones}

La Ciudad de México enfrenta un cáncer sistémico en más de una área; el transporte no es un factor, 
es el resultado de múltiples esfuerzos de diferentes actores. Si este es ineficiente, el problema no 
se debe a la absurda idea de \termino{Meter más líneas de camiones} o 
\termino{Desplazar más cantidad de gente} a un cierto punto. Obedece a un resultado multifactorial: 
\begin{itemize}
    \item \textbf{Económico:} El transporte público debe ser accesible a todas las personas sin segregar a nadie por su
            estrato socio-económico. Debe fijarse una tarifa accesible y regulada por un periodo de tiempo prolongado
            que permita darle un mantenimiento adecuado a sus instalaciones, infraestructura y sobre todo, a sus trabajadores,
            sindicatos y directivos. No se puede tener un transporte de calidad sin darle una calidad digna a quienes lo operan.

    \item \textbf{Social:} La movilidad no debe verse como un castigo sistémico que te recuerda que \termino{eres pobre}. 
            Países del primer mundo como la ciudad de Helkensinki, Japón y el tren de alta velocidad en España revelan un cambio 
            de paradigma; el transporte público no debe segregar, debe unir y ser un espacio seguro, transitable e incluyente para 
            cualquier género, preferencia sexual, estrato socio-económico y persona con capacidades diferentes; 
            tanto psicológico como físico.

    \item \textbf{Geométrico - Matemático:} Quizás el menos explorado de la lista. Existen infinidad de artículos que hablan sobre la 
            ineficacia de la red de transporte en cualquier región del mundo atacándola con más infraestructura, recursos o deterioro social. 
            Pero pocos son las voces que tratan el fallo como un problema de logística y planeación. Incluso el modelo DOT: Desarrollo Orientado 
            al Transporte no contempla cambios en la distribución de su geometría; es decir pasar de un paradigma radial, donde todas las rutas 
            convergen en el centro, a un modelo circular, donde el centro se vea intacto y se pueda llegar a otra región de la urbe sin saturar 
            el sistema en las llamadas \termino{Horas pico}.

    \item \textbf{Cultural:} Muchos son los autores que describen a México como un país rico en cultura, tradiciones y calidez humana. Esto es 
            totalmente cierto. Pero como lo han revelado figuras controversiales como lo son \textbf{Octavio Paz} y \textbf{Gustavo Díaz Ordaz}, 
            uno de los mayores problemas de México, es el México mismo. Su población esta concentrada en sobrevivir 
            (esto debido a la precarización que el sistema les ha obligado a vivir), pero incluso cuando la oportunidad de cuidar el ambiente o las instalaciones 
            de una obra pública (como lo es el mismo Metro de la Ciudad de México), se cometen actos de bandalísmo por parte de la misma población. 
            Existe una cultura impregnada en la población actual que se define en el refrán popular: \termino{El que no tranza, no avanza} y 
            \termino{De que lloren en tu casa a que lloren en la mía, mejor que lloren en la tuya}.

    \item \textbf{Ambiental:} Quizás el punto más avanzado en términos generacionales, se ha despertado una conciencia en torno a los daños ambientales que 
            puede causar el uso prolongado del automóvil y la quema de hidrocarburos al medio ambiente. El urbanismo ha decidido atacar este punto 
            con el término de \termino{Desarrollo Sustentable}. Sin embargo y como se exploró en secciones anteriores de este ensayo, casi ningún proyecto actual 
            en México es \termino{Sustentable} ni \termino{Sostenible}, ya que detrás de la construcción de toda obra pública en los últimos 20 años, siempre 
            ha habido intereses de empresas privadas, corrupción o nepotismo. A pesar de ello, el transporte público es la única vía para hacer una 
            movilidad más amigable con el ambiente, no solo en sus emisiones de carbono, si no en términos económicos, se vuelve más rentable a lo largo del tiempo.

    \item \textbf{Tecnológico:} En el año que se escribe este ensayo (2026), se tiene un mayor acceso al hardware (elementos físicos de una computadora) y 
            un mayor conocimiento del software (elementos internos de una tecnología como lo son programas o lenguajes de programación). Muchas de estas herramientas 
            son de uso libre y gratuito. Cualquier persona con interés, tiempo y dedicación puede construir su propio sistema usando una computadora básica, 
            manuales de código extraídos de internet y metodologías o investigaciones propias tanto profundas como superficiales. Sin embargo, el avance de esta 
            rama se ha visto limitado por la cultura del entretenimiento y el consumo. Los celulares de hoy día pueden ejecutar videojuegos que hace 20 años necesitaban 
            un servidor dedicado y las mejoras en ellos son (casi siempre) en el apartado gráfico. Las computadoras de escritorio hoy día son más poderosas que cualquier 
            herramienta de cálculo en la antigüedad, y sin embargo, se sigue avanzando en la cultura del entretenimiento. Las pocas personas que tiene acceso a cambiar este 
            pensamiento yacen con el poder político o económico en sus manos y muchas de ellas no lo usan en el beneficio de la mayoría, es por esto que el software privativo 
            o empresarial ha experimentado una demanda progresiva a lo largo de los últimos 30 años. A su vez, el entendimiento de mejora tecnológica casi siempre se ve 
            con la adopción de \termino{Tecnologías inovadoras} en ves de tecnologías útiles. El ejemplo más claro es los tramites administrativos, tener digitalizados todos 
            los expedientes judiciales de una persona durante un proceso penal no sirve de nada si estos aún deben pasar por la burocracia sistemática del estado. 
            No se tiene a la fecha del día de hoy, un software de uso libre que proporcione una herramienta real para el análisis del transporte público y existen carencias latentes 
            en el sistema tecnológico mexicano; si se quiere consultar una base de datos del sistema de transporte público de la Ciudad de México, tendrás que enfrentarte 
            con cientos de miles de archivos difíciles de entender para cualquier ser humano. Sin embargo, la tesis 
            \textbf{\termino{Transportes Anillares y su importancia en la Ciudad de México y el área metropolitana}} (ver sección de Releases) \cite{galigaribaldi_tesis}, 
            resuelve este conflicto creando dos herramientas de uso libre: 
            \textbf{Apimetro: La primera API del sistema de transporte público de la Ciudad de México} (Ver Página de inicio apimetro.dev) \cite{galigaribaldi_apimetro}y 
            \textbf{VFTModel: La primera herramienta de software libre que transforma datos espaciales a grafos matemáticos para su posterior análisis} \cite{galigaribaldi_vftmodel}. 
\end{itemize}